% =====================================================================
%  Section 8 (opening) and 8.1--8.3: Cole--Hopf, exact solution,
%  vanishing viscosity.  8.4 and 8.5 are in ext_godunov.tex and
%  ext_nwave.tex and can be removed independently from main.tex.
% =====================================================================

\section{Extensions}
\label{sec:extensions}

The core theory (Sections~\ref{sec:characteristics}--\ref{sec:riemann-problem})
is complete without this section. Subsections~\ref{subsec:cole-hopf}--%
\ref{subsec:vanishing-viscosity} reintroduce viscosity and connect the
viscous and inviscid theories;
Subsection~\ref{subsec:arctan-complete} completes the worked example of
Sections~\ref{sec:characteristics}--\ref{sec:shock-formation} using them;
Subsections~\ref{subsec:godunov} and \ref{subsec:n-wave} are self-contained
applications of the Riemann solver and of the entropy theory.

\subsection{The Cole--Hopf transformation}
\label{subsec:cole-hopf}

(Due to Cole \cite{cole1951} and Hopf \cite{hopf1950}.)

This extension solves the viscous Burgers' equation
\eqref{eq:viscous-burgers} exactly by linearising it to the heat equation.
The key observation is that, in one dimension, any $u$ is the
$x$-derivative of a potential, namely $-2\nu\,\partial_x$ of a logarithm.

\begin{theorem}[Cole--Hopf transformation]
\label{thm:cole-hopf}
Let $u$ be a classical solution of $u_t+u\,u_x=\nu\,u_{xx}$ on
$\mathbb{R}\times[0,T)$, and set
\begin{equation}
  \label{eq:cole-hopf-def}
  \hat\psi(x,t)=-\frac{1}{2\nu}\int_0^x u(y,t)\,dy,
  \qquad
  c(t)=\int_0^t\Big(\frac{u(0,\tau)^2}{4\nu}-\frac{u_x(0,\tau)}{2}\Big)d\tau,
  \qquad
  \theta(x,t)=\exp\big(\hat\psi(x,t)+c(t)\big)>0 .
\end{equation}
Then $\theta$ solves the heat equation $\theta_t=\nu\,\theta_{xx}$, and
$u=-2\nu\,\theta_x/\theta$.
\end{theorem}

\begin{proof}
The identity $u=-2\nu\theta_x/\theta$ is immediate: $\theta_x/\theta=
\hat\psi_x=-u/(2\nu)$, since $c$ depends on $t$ only. For the heat
equation, let $w=\hat\psi_t-\nu\hat\psi_{xx}-\nu\hat\psi_x^2$. Using
$\hat\psi_x=-u/(2\nu)$, $\hat\psi_{xx}=-u_x/(2\nu)$,
$\hat\psi_{xt}=-u_t/(2\nu)$ and $\hat\psi_{xxx}=-u_{xx}/(2\nu)$,
\begin{equation}
  w_x=\hat\psi_{xt}-\nu\hat\psi_{xxx}-2\nu\hat\psi_x\hat\psi_{xx}
  =-\frac{1}{2\nu}\big(u_t+u\,u_x-\nu\,u_{xx}\big)=0 .
\end{equation}
So $w$ is a function of $t$ alone, and may be evaluated at $x=0$, where
$\hat\psi(0,t)=0$ gives $\hat\psi_t(0,t)=0$; hence
\begin{equation}
  w(t)=w(0,t)=\frac{u_x(0,t)}{2}-\frac{u(0,t)^2}{4\nu}=:C(t).
\end{equation}
Put $\psi=\hat\psi+c$, so $c'=-C$ and $\psi_t-\nu\psi_{xx}-\nu\psi_x^2=
C+c'=0$. Since $\theta=e^{\psi}$ gives $\theta_t=\theta\psi_t$ and
$\theta_{xx}=\theta(\psi_{xx}+\psi_x^2)$, this is $\theta_t=\nu\theta_{xx}$.
\end{proof}

\begin{remark}
The correction $c(t)$ cannot be omitted in general. The naive function
$e^{\hat\psi}=\exp\big(-\tfrac1{2\nu}\int_0^x u\,dy\big)$, normalised to
equal $1$ at $x=0$, satisfies $\partial_t\theta=\nu\theta_{xx}+C(t)\theta$,
which is the heat equation only if $C\equiv0$. The factor $e^{c(t)}$ is
independent of $x$, so it changes neither $u=-2\nu\theta_x/\theta$ nor the
initial data ($c(0)=0$).
\end{remark}

\subsection{Exact solution via the heat kernel}
\label{subsec:cole-hopf-solution}

By Theorem~\ref{thm:cole-hopf}, solving \eqref{eq:viscous-burgers} reduces
to solving the heat equation for $\theta$ with initial data
\begin{equation}
  \label{eq:theta0}
  \theta_0(x) = \exp\left(-\frac{1}{2\nu}\int_0^x u_0(y)\, dy\right),
\end{equation}
which, for bounded $u_0$, grows at most like $e^{\|u_0\|_\infty|x|/(2\nu)}$.
The Cauchy problem for the heat equation has a unique solution in this
growth class (see \cite[\S2.3.3]{evans2010}), given by convolution with
the heat kernel,
\begin{equation}
  \label{eq:theta-heat-kernel}
  \theta(x,t) = \int_{\mathbb{R}} \frac{1}{\sqrt{4\pi\nu t}}\,
    \exp\left(-\frac{(x-y)^2}{4\nu t}\right) \theta_0(y)\, dy.
\end{equation}
Substituting \eqref{eq:theta0}--\eqref{eq:theta-heat-kernel} into
$u=-2\nu\,\theta_x/\theta$ and differentiating under the integral gives
the explicit \emph{Cole--Hopf formula}
\begin{equation}
  \label{eq:cole-hopf-formula}
  u(x,t) = \frac{\displaystyle\int_{\mathbb{R}} \frac{x-y}{t}\, \exp\left(-\frac{G(x,y,t)}{2\nu}\right) dy}
                {\displaystyle\int_{\mathbb{R}} \exp\left(-\frac{G(x,y,t)}{2\nu}\right) dy},
  \qquad
  G(x,y,t) = \int_0^y u_0(s)\, ds + \frac{(x-y)^2}{2t}.
\end{equation}

\begin{corollary}[Global smoothness; Hopf \cite{hopf1950}]
\label{cor:global-smoothness}
For every bounded measurable $u_0$ and every $\nu>0$, formula
\eqref{eq:cole-hopf-formula} defines a solution of
\eqref{eq:viscous-burgers} that is $C^\infty$ for all $t>0$; it is the
unique bounded solution with data $u_0$.
\end{corollary}

\noindent That \eqref{eq:cole-hopf-formula} solves the equation follows by
reversing the computation of Theorem~\ref{thm:cole-hopf}: if $\theta>0$
solves the heat equation, then $u=-2\nu\theta_x/\theta$ solves
\eqref{eq:viscous-burgers}, and $\theta$ given by
\eqref{eq:theta-heat-kernel} is smooth and positive for $t>0$. The
statement that the initial data are attained and that the bounded
solution is unique for merely measurable $u_0$ is due to Hopf
\cite{hopf1950}; it is cited, not reproved here.

Unlike the inviscid equation (Proposition~\ref{prop:shock-time}),
solutions of \eqref{eq:viscous-burgers} stay smooth for all time, for
any $\nu>0$ however small: advection and viscous diffusion together do
not by themselves produce singularities. This is a statement about Burgers' equation, not about
Navier--Stokes: Burgers' scalar, one-dimensional structure (it obeys a
maximum principle and has no pressure or vortex stretching) means that its
regularisation mechanism is quite different from the open regularity
question in three dimensions.

\subsection{The vanishing viscosity limit}
\label{subsec:vanishing-viscosity}

The Cole--Hopf solution also explains why the entropy conditions of
Section~\ref{sec:entropy-condition} are the physically correct selection
principle, rather than an arbitrary one.

\begin{theorem}[Vanishing viscosity limit; Hopf, 1950]
\label{thm:vanishing-viscosity}
Let $u_0\in L^\infty(\mathbb{R})$, let $u^\nu$ denote the solution of
\eqref{eq:viscous-burgers} with data $u_0$ given by
\eqref{eq:cole-hopf-formula}, and let $u$ denote the entropy solution of
\eqref{eq:inviscid-burgers} with the same data, in the sense of
Definition~\ref{def:entropy-solution} (for Riemann data this is
Theorem~\ref{thm:riemann-solution}). Then $u^\nu \to u$ pointwise almost
everywhere, and in $L^1_{\mathrm{loc}}$, as $\nu \to 0^+$ \cite{hopf1950}.
\end{theorem}

The proof is a Laplace-method asymptotic analysis of
\eqref{eq:cole-hopf-formula} as $\nu\to0$, and is not reproduced here. Its
mechanism, however, is worth recording, because it is used again in
Section~\ref{subsec:n-wave}.

\begin{remark}[The Lax--Oleinik formula]
\label{rem:lax-oleinik}
As $\nu\to0$ the weight $\exp(-G/2\nu)$ in
\eqref{eq:cole-hopf-formula} concentrates at the minimisers of
$y\mapsto G(x,y,t)$. At every $(x,t)$ where the minimiser
$y^*(x,t)$ is unique, the limit is
\begin{equation}
  \label{eq:lax-oleinik}
  u(x,t)=\frac{x-y^*(x,t)}{t},\qquad
  y^*(x,t)=\operatorname*{arg\,min}_y\Big\{V(y)+\frac{(x-y)^2}{2t}\Big\},
  \quad V(y)=\int_0^y u_0 ,
\end{equation}
the \emph{Lax--Oleinik formula} \cite{lax1957,oleinik1957}, \cite[\S3.4.2]{evans2010}. A jump occurs
where two minimisers $y_1<y_2$ tie, with $u_\ell=(x-y_1)/t>u_r=(x-y_2)/t$:
the viscous limit produces only \emph{downward} jumps, which is exactly
the content of the Lax condition \eqref{eq:lax-burgers}. More
generally, $x\mapsto y^*(x,t)$ is nondecreasing (a standard property of the
minimiser, \cite[\S3.4.2]{evans2010}), so for $z>0$
\begin{equation*}
  u(x+z,t)-u(x,t)=\frac{z-\big(y^*(x+z,t)-y^*(x,t)\big)}{t}\ \le\ \frac zt,
\end{equation*}
which is the sharp Oleinik bound \eqref{eq:oleinik}. Together with the fact
that \eqref{eq:lax-oleinik} defines a weak solution
\cite[\S3.4.2]{evans2010}, this shows that the Lax--Oleinik solution is an
entropy solution in the sense of Definition~\ref{def:entropy-solution}.
\end{remark}

\subsubsection{Viscous shock profiles and the Lax condition}
\label{subsubsec:viscous-profiles}

The vanishing-viscosity picture can be made quite explicit at a single jump.
It also shows directly that the Lax condition is not an arbitrary rule.

\begin{proposition}[Viscous shock profiles]
\label{prop:viscous-profile}
Let $u_\ell\neq u_r$ and $\nu>0$. The viscous Burgers' equation
\eqref{eq:viscous-burgers} has a travelling-wave solution $u(x,t)=\phi(x-st)$,
with $\phi\in C^2$ bounded and $\phi(\xi)\to u_\ell$, $\phi(\xi)\to u_r$ as
$\xi\to-\infty$, $+\infty$ (and $\phi'\to0$), \emph{if and only if}
$u_\ell>u_r$. In that case $s=(u_\ell+u_r)/2$ and, with $\Delta=u_\ell-u_r$,
\begin{equation}
  \label{eq:viscous-profile}
  \phi(\xi)=s-\frac{\Delta}{2}\tanh\!\Big(\frac{\Delta\,\xi}{4\nu}\Big).
\end{equation}
Moreover $\|\phi-(s-\tfrac\Delta2\operatorname{sgn}\xi)\|_{L^1(\mathbb{R})}
=4\ln2\;\nu$, independent of $s$ and $\Delta$.
\end{proposition}

\begin{proof}
Substituting $u=\phi(x-st)$ gives $-s\phi'+\phi\phi'=\nu\phi''$, and
integrating once from $-\infty$ (where $\phi'\to0$),
\begin{equation*}
  \nu\phi'=\tfrac12\phi^2-s\phi-C,\qquad C=\tfrac12u_\ell^2-su_\ell .
\end{equation*}
Letting $\xi\to+\infty$ shows the right-hand side must also vanish at $u_r$,
so $u_\ell,u_r$ are the two roots of the quadratic: $s=(u_\ell+u_r)/2$
(the Rankine--Hugoniot speed) and
\begin{equation}
  \label{eq:profile-ode}
  \nu\phi'=\tfrac12(\phi-u_\ell)(\phi-u_r).
\end{equation}
This is a first-order autonomous ODE, so $\phi$ cannot cross the equilibria
$u_\ell,u_r$ and, to connect them, must lie strictly between them. If
$u_\ell<u_r$, then for $\phi\in(u_\ell,u_r)$ the right side of
\eqref{eq:profile-ode} is \emph{negative}, so $\phi$ decreases, yet it must
rise from $u_\ell$ to $u_r$: no profile exists. If $u_\ell>u_r$ the right side
is negative on $(u_r,u_\ell)$, consistent with a decrease from $u_\ell$ to
$u_r$, and \eqref{eq:profile-ode} is solved by \eqref{eq:viscous-profile}.
For the $L^1$ distance, by oddness,
\begin{equation*}
  \int_{\mathbb{R}}\tfrac\Delta2\Big|\tanh\tfrac{\Delta\xi}{4\nu}-\operatorname{sgn}\xi\Big|\,d\xi
  =\Delta\int_0^\infty\Big(1-\tanh\tfrac{\Delta\xi}{4\nu}\Big)d\xi
  =\Delta\cdot\frac{4\nu}{\Delta}\ln2=4\ln2\,\nu,
\end{equation*}
using $\int_0^\infty(1-\tanh s)\,ds=\ln2$.
\end{proof}

\begin{remark}
(i) \emph{The Lax condition from viscosity.} A jump with $u_\ell<u_r$ is not
the limit of viscous travelling waves, whereas every jump with $u_\ell>u_r$ is,
with thickness $\sim\nu/\Delta\to0$. This derives the admissibility condition
\eqref{eq:lax-burgers} from the vanishing-viscosity principle directly, and
explains why $u_A$ of Section~\ref{sec:nonuniqueness} is rejected: it contains
an upward jump, which no vanishing-viscosity limit can produce.
(ii) \emph{Galilean invariance.} Equation \eqref{eq:viscous-burgers} is
invariant under $u\mapsto u+c$, $x\mapsto x+ct$, which maps a profile of
speed $s$ to one of speed $s-c$ and leaves $L^1$ distances between $u$-values
unchanged. This is why the constant $4\ln2\,\nu$ does not depend on the speed
$s$: for the stationary shock of strength $a$ ($s=0$, $\Delta=2a$) it is the
profile $-a\tanh(ax/2\nu)$ used in Section~\ref{subsec:arctan-complete}, and
for the shock $1\to0$ of Section~\ref{subsubsec:numerical-viscosity}
($s=\tfrac12$, $\Delta=1$) it is the same profile boosted by $\tfrac12$.
\end{remark}
